\documentclass[11pt,a4paper]{article}
\usepackage{curso}
\begin{document}
\portada{Lección}{0.3}{Repaso de probabilidad}

\noindent\textbf{Unidad 0.}\quad\textbf{Objetivos:} manejar Bernoulli, categóricas y gaussianas; transformar
gaussianas multivariadas; calcular entropía de Shannon y divergencia KL. \textbf{Prerrequisitos:} cálculo,
álgebra matricial básica.

\section{¿Por qué importa?}
Un computador termodinámico devuelve \emph{muestras}. Para saber qué calcula (la distribución), con qué
precisión (sus momentos y el error de promediar) y cuánta información contiene (entropía), necesitamos el
lenguaje de la probabilidad.

\section{Variables discretas}
\begin{definicion}[ (Bernoulli y categórica)]
$x\sim\mathrm{Bernoulli}(p)$ si $P(x=1)=p$, $P(x=0)=1-p$; $\E[x]=p$, $\Var[x]=p(1-p)$.
$x\sim\mathrm{Cat}(\pi_1,\dots,\pi_K)$ si $P(x=k)=\pi_k$ con $\pi_k\ge0$ y $\sum_k\pi_k=1$.
\end{definicion}
\textbf{Ley de los grandes números.} Si $x_1,\dots,x_n$ son i.i.d. con media $\mu$ y varianza $\sigma^2$, la
media muestral $\bar x_n$ converge a $\mu$ y su error estándar es $\sigma/\sqrt n$. Para Bernoulli($p$), el
error estándar de $\hat p$ es $\sqrt{p(1-p)/n}$.

En el curso el espín de un pbit toma valores $s\in\{-1,+1\}$; si $x=(s+1)/2$ es Bernoulli($p$), entonces
$\E[s]=2p-1$ (Unidad 4).

\section{La gaussiana y el teorema central del límite}
\begin{definicion}[ (gaussiana univariada)]
$p(x)=\frac{1}{\sqrt{2\pi\sigma^2}}\exp\!\big(-\frac{(x-\mu)^2}{2\sigma^2}\big)$, con $\E[x]=\mu$, $\Var[x]=\sigma^2$.
\end{definicion}
\begin{resultado}[ (TCL)]
Si $x_i$ son i.i.d. con media $\mu$ y varianza $\sigma^2<\infty$, $\frac{\sum_{i=1}^n x_i-n\mu}{\sigma\sqrt n}\to\N(0,1)$
en distribución. En la animación, la suma de 20 Bernoulli(1/2) (binomial) se aproxima por $\N(10,5)$.
\end{resultado}
El ruido térmico de un resistor o la fuerza aleatoria sobre una partícula browniana son sumas de enormes
números de choques microscópicos casi independientes: por eso se modelan como gaussianos (Unidad 2).

\section{Gaussiana multivariada}
$x\sim\N(\mu,\Sigma)$ en $\R^d$ con $\Sigma$ simétrica definida positiva (SPD):
\[
p(x)=\frac{1}{\sqrt{\det(2\pi\Sigma)}}\exp\!\Big(-\tfrac12(x-\mu)^\top\Sigma^{-1}(x-\mu)\Big).
\]
\begin{proposicion}
Si $y\sim\N(0,I)$ y $x=Ay+b$, entonces $x\sim\N(b,AA^\top)$.
\end{proposicion}
\begin{proof}
Linealidad: $\E[x]=A\E[y]+b=b$. $\mathrm{Cov}(x)=\E[(Ay)(Ay)^\top]=A\,\E[yy^\top]A^\top=AA^\top$. Una
transformación afín de una gaussiana es gaussiana (su función característica sigue siendo exponencial de
una cuadrática), y una gaussiana queda determinada por media y covarianza.
\end{proof}
Geométricamente, las curvas de nivel son elipses $\{x:(x-\mu)^\top\Sigma^{-1}(x-\mu)=c\}$ con ejes en los
vectores propios $v_i$ de $\Sigma$ y semiejes $\propto\sqrt{\lambda_i}$ (lección 0.4). Recíprocamente, para
muestrear $\N(\mu,\Sigma)$ basta $x=\mu+Ly$ con $\Sigma=LL^\top$ (Cholesky).

\section{Momentos}
Media $\mu=\E[x]$; covarianza $\Sigma=\E[(x-\mu)(x-\mu)^\top]$; momentos de orden superior
$\E[x^k]$. Para gaussianas, los dos primeros momentos determinan todo. \textbf{Anticipo (Unidades 2 y 6):}
el proceso de Ornstein--Uhlenbeck $\dd x=-(Ax-b)\dd t+\sqrt{2\beta^{-1}}\dd W$ tiene estado estacionario
$\N(A^{-1}b,\beta^{-1}A^{-1})$: la media resuelve $Ax=b$ y la covarianza da $A^{-1}$.

\section{Entropía de Shannon}
\begin{definicion}[ (entropía)]
$H(x)=-\sum_ip_i\log_2p_i$ bits; con $\ln$ se mide en nats ($1$ nat $=1/\ln2\approx1{,}443$ bits).
Para Bernoulli: $h(p)=-p\log_2p-(1-p)\log_2(1-p)$, máxima ($1$ bit) en $p=1/2$.
\end{definicion}
Interpretación: $-\log_2p_i$ es la sorpresa del resultado $i$; $H$ es la sorpresa promedio y el número
mínimo promedio de bits por símbolo para codificar (teorema de codificación de fuente).
Para densidades se define la entropía diferencial $h=-\int p\ln p$; para $\N(\mu,\Sigma)$ en $\R^d$:
\[h=\tfrac12\ln\det(2\pi e\,\Sigma)=\tfrac d2\ln(2\pi e)+\tfrac12\ln\det\Sigma,\]
que puede ser negativa (si $\Sigma$ es pequeña).

\section{Divergencia de Kullback--Leibler}
\begin{definicion}[ (KL)]
$\KL(p\|q)=\E_p\big[\ln\frac{p}{q}\big]$. Por la desigualdad de Jensen (Gibbs), $\KL(p\|q)\ge0$ con igualdad
si y solo si $p=q$ casi en todas partes. No es simétrica ni cumple la desigualdad triangular.
\end{definicion}
Para gaussianas univariadas:
\[\KL(\N(\mu_1,\sigma_1^2)\|\N(\mu_2,\sigma_2^2))=\ln\frac{\sigma_2}{\sigma_1}+\frac{\sigma_1^2+(\mu_1-\mu_2)^2}{2\sigma_2^2}-\frac12.\]
\begin{ejemplo}
$h(0{,}9)=-0{,}9\log_20{,}9-0{,}1\log_20{,}1=0{,}137+0{,}332=0{,}469$ bits.
$\KL(\N(0,1)\|\N(1,4))=\ln2+\frac{1+1}{8}-\frac12\approx0{,}443$ nats. Con $\N(0,4)$ en lugar de $\N(1,4)$:
$\KL(\N(0,1)\|\N(0,4))=0{,}318$ y $\KL(\N(0,4)\|\N(0,1))=0{,}807$: la asimetría es evidente.
\end{ejemplo}

\begin{aplicacion}
Pbit $=$ Bernoulli controlada por un campo (Unidad 4). Estado estacionario del OU $=$ gaussiana
multivariada (Unidades 2, 6). Entropía $\to$ costo de Landauer (Unidad 1). KL $\to$ aprendizaje de máquinas
de Boltzmann, ELBO de difusión, geometría de la información y gradiente natural (Unidades 4, 8, 9).
\end{aplicacion}

\begin{error}
\begin{itemize}[nosep]
\item $\mathrm{Cov}(Ay+b)=A\,\mathrm{Cov}(y)A^\top$, no $A^2\mathrm{Cov}(y)$.
\item KL no es distancia. Elegir $\KL(p\|q)$ o $\KL(q\|p)$ cambia el comportamiento (cubrir modos frente a buscar modos).
\item Bits ($\log_2$) frente a nats ($\ln$).
\item La entropía diferencial puede ser negativa y no es invariante ante cambios de escala.
\end{itemize}
\end{error}

\section{Resumen}
Bernoulli/categórica/gaussiana; LGN y TCL; $x=Ay+b\Rightarrow\N(b,AA^\top)$; $H=-\sum p\log p$; $\KL\ge0$ no simétrica.

\begin{thebibliography}{9}\small
\bibitem[3]{levin} D. A. Levin, Y. Peres y E. L. Wilmer, \emph{Markov Chains and Mixing Times}, AMS, 2009.
\bibitem[46]{holderrieth} P. Holderrieth y E. Erives, An introduction to flow matching and diffusion models, 2025. \url{https://arxiv.org/abs/2506.02070}
\bibitem[c]{cover} (Complementaria, fuera del pensum) T. M. Cover y J. A. Thomas, \emph{Elements of Information Theory}, 2.ª ed., Wiley, 2006.
\end{thebibliography}
\end{document}
